\documentclass[12pt]{article}
\usepackage[utf8]{inputenc}
\usepackage[T1]{fontenc}
\usepackage[margin=1in]{geometry}
\usepackage{setspace}
\usepackage{xcolor}
\usepackage[colorlinks=true, urlcolor=blue!60!black, linkcolor=blue!60!black]{hyperref}
\setlength{\parindent}{0em}
\setlength{\parskip}{0.6em}
\renewcommand{\baselinestretch}{1.0}
\begin{document}
\begin{center}
{\LARGE \textit{Sveen v. Melin}, 584 U.S. 811 (2018)\par}
\vspace{0.5em}
{\large Case Brief --- Two-Page Edition\par}
\end{center}

\section{Memory Jogger}
A retroactive revocation-on-divorce statute that removed an ex-wife as life-insurance beneficiary did not substantially impair the contract, so it did not violate the Contracts Clause (\href{https://supreme.justia.com/cases/federal/us/584/16-1432/}{584 U.S.\ 811}).

\section{Detailed Case Facts}
Mark Sveen named his wife, Kaye Melin, as beneficiary of his life-insurance policy in 1998, with his children as contingent beneficiaries. In 2002 Minnesota extended its revocation-on-divorce rule to life insurance. The couple divorced in 2007 without addressing the policy, and Mark never changed the designation. After his death in 2011, Melin and the children both claimed the proceeds.

\section{Procedural History}
In MetLife's interpleader action, the district court ruled for the children; the Eighth Circuit reversed under the Contracts Clause.

\section{Judicial Votes}
Reversed, 8--1. Kagan wrote for Roberts, Kennedy, Thomas, Ginsburg, Breyer, Alito, and Sotomayor; Gorsuch dissented.

\section{Holding}
Under the two-step Contracts Clause framework, the statute fails the threshold requirement of ``substantial impairment,'' so the Court did not reach means-ends review.

\section{Analysis of Opinions}
Kagan found three features decisive: the law reflects likely intent, since ``the average Joe does not want his ex inheriting what he leaves behind''; it does no more than a divorce court could always do; and the policyholder can undo the default ``with the stroke of a pen,'' like the retroactive recording and filing laws upheld since the 1800s, including \textit{Texaco v.\ Short}. Gorsuch, dissenting, argued that the Clause forbids any retroactive impairment, that naming the beneficiary is the ``whole point'' of the contract, and that the default traps inattentive policyholders.

\section{Examples: Future Applications}
\textbf{Same-Side:} a State may apply its divorce-revocation rule to existing bank accounts payable on death. \textbf{Opposite-Side:} retroactively voiding all non-competes with no opt-out substantially impairs contracts. \textbf{Fence-Sitter:} a revocation rule that can be undone only by court petition.

\section{Critique}
The ``stroke of a pen'' rationale assumes policyholders know the law changed, and folding reasonableness into step one lets courts skip tailoring review.

\section{Key Quotations}
\begin{quote}
``In answering that question, the Court has considered the extent to which the law undermines the contractual bargain, interferes with a party's reasonable expectations, and prevents the party from safeguarding or reinstating his rights.''
\end{quote}
\noindent --- Justice Kagan, for the Court, 584 U.S.\ at 819.

\end{document}
